\documentclass[10pt,a4paper]{amsart}
\usepackage[margin=19mm]{geometry}
\usepackage{amsmath,amssymb,mathtools}
\usepackage[scr=rsfs,cal=euler]{mathalfa}
\usepackage{xfrac}
\usepackage[unicode,hidelinks,bookmarks=false]{hyperref}
\newcommand{\ie}{i.e.\ }
\newcommand{\cf}{cf.\ }
\newcommand{\resp}{respectively\ }
\newcommand{\Subsection}{\subsection}
\providecommand{\nobreakdash}{\nobreak\hbox{-}}
\setlength{\emergencystretch}{2em}
\multlinegap0pt
\usepackage{dsfont}
\usepackage{amssymb}
\usepackage[shortlabels]{enumitem}
\usepackage{float}
\usepackage{xcolor}
\newtheorem{theorem}{Theorem}
\title{From subtractive ideals of semirings to deductive and inductive sets in general algebras}
\author{Elena Caviglia, Amartya Goswami, Zurab Janelidze, Luca Mesiti,, Vaino T. Shaumbwa}
\date{Source: arXiv:2602.01178v2 (CC BY 4.0)}
\begin{document}
\maketitle

\noindent\emph{Source and licence.} From subtractive ideals of semirings to deductive and inductive sets in general algebras. Version arXiv:2602.01178v2 (CC BY 4.0), distributed by arXiv under the Creative Commons Attribution 4.0 International licence, http://creativecommons.org/licenses/by/4.0/, as recorded in the arXiv licence metadata for this exact version and shown on the abstract page https://arxiv.org/abs/2602.01178v2. Full licence notice: \texttt{licenses/arxiv-2602.01178-CC-BY-4.0.txt}.\par\smallskip

\noindent\textit{Editorial scope.} Counted unit: the article's theorem DetThmLatticeUniversalRank (source0.tex lines 2492--2501). The article's complete original proof of it is delivered with it (source0.tex lines 2503--2622, one \texttt{proof} environment of 120 lines). No mathematics is written, repaired or re-derived: the statement and the proof are the article's own lines, in the article's own notation, and no retained line is substituted or re-flowed.  No further source-local context is needed: every cross-reference of every retained line resolves inside the two delivered spans.  Nothing was omitted from any retained span, so the package carries no omission record.  No retained line contains a citation, so the wrapper carries no bibliography and no bibliography entry.  No retained line contains a figure environment, an image-inclusion command or drawing code, so no image is delivered and the package declares no asset dependency.  Editorial formatting only, in addition: an AMS \texttt{amsart} preamble that restates the article's own declarations and macros used by the retained lines, namely the environments \texttt{theorem} on separate counters, together with the standard packages those lines need.  The retained environments print in this document as Theorem 1 in order of appearance, whereas the article prints the same items with the numbering of its own full document.  The complete native source of the article as distributed in the arXiv e-print for this exact version is bundled unchanged as \texttt{sources/193769/source0.tex}.
\begin{theorem}\label{DetThmLatticeUniversalRank}
Let $V$ be a variety of lattices, possibly with named constants. Then
\[
r_{\mathcal I}^{\mathrm{univ}}(V)\leqslant1
\quad\Longleftrightarrow\quad
V\text{ is distributive}.
\]
Every non-trivial such distributive variety has
$\operatorname{urk}(V)=(1,1)$.
\end{theorem}

\begin{proof}
First let $L$ be distributive, fix $a\in L$, and let $I\subseteq L$
be non-empty. Put
\[
x\preceq_a y\quad\Longleftrightarrow\quad
x\wedge a\leqslant y\leqslant x\vee a,
\qquad
\theta_I=\operatorname{Cg}_L(I\times\{a\}).
\]
Distributivity shows that $\preceq_a$ is reflexive, transitive and
compatible. We claim that the semicongruence generated by
$I\times\{a\}$ is
\begin{equation}\label{DetEqDistributiveStar}
R_{I,a}=\theta_I\cap\preceq_a.
\end{equation}

For finite $I$, put $\ell=a\wedge\bigwedge I$ and
$h=a\vee\bigvee I$. Both $(\ell,a)$ and $(h,a)$ belong to $R_{I,a}$,
and the polynomial $(s\vee U)\wedge V$ recovers $(s,a)$ from these two
pairs, for each $s\in I$. Thus they generate $R_{I,a}$, and
$\theta_I=\operatorname{Cg}_L(\ell,h)$. Moreover,
\begin{equation}\label{DetEqDistributivePrincipalCongruence}
x\mathrel{\theta_I}y
\quad\Longleftrightarrow\quad
x\wedge\ell=y\wedge\ell\ \text{ and }\ x\vee h=y\vee h.
\end{equation}
To verify this, the right-hand side is the kernel of a lattice
homomorphism and contains $(\ell,h)$. Conversely, in the quotient by
$\operatorname{Cg}_L(\ell,h)$, denote the common image of $\ell,h$ by
$c$. Equality of both meets and both joins with $c$ forces the two
images to be equal, since in a distributive lattice
\[
x=x\wedge(y\vee c)
=(x\wedge y)\vee(x\wedge c)
=(x\wedge y)\vee(y\wedge c)=y.
\]
Now $R_{I,a}\subseteq\theta_I\cap\preceq_a$ follows by compatibility.
For the reverse inclusion, suppose $x\mathrel{\theta_I}y$ and
$x\preceq_a y$. The polynomial
\[
p(U,V)=(x\wedge V)\vee(y\wedge U)\vee(x\wedge y)
\]
satisfies $p(\ell,h)=x$ by
\eqref{DetEqDistributivePrincipalCongruence}, and $p(a,a)=y$ by the
inequalities defining $\preceq_a$. This proves
\eqref{DetEqDistributiveStar} for finite $I$. The general case follows
by taking directed unions over finite non-empty subsets of $I$.

Write $R=R_{I,a}$. Both $RI$ and $IR$ contain $I$ and lie in the
$\theta_I$-class of $a$. Replacing $I$ by either set therefore leaves
$\theta_I$ unchanged, and \eqref{DetEqDistributiveStar} shows that it
also leaves $R$ unchanged. Transitivity now gives
\[
\mathcal I_a^2(I)=R^2I=RI,
\qquad
\mathcal D_a^2(I)=IR^2=IR.
\]

For the converse, let $L\in V$ be non-distributive. Its order relation
\[
B=\{(u,v)\in L^2\mid u\leqslant v\}
\]
is a subalgebra of $L^2$, including any named constants. We use the
following one-step obstruction. If $\alpha=(d,e)\in B$ and
$T\subseteq B$ satisfies $d\leqslant\pi_2(t)$ for every $t\in T$,
then
\begin{equation}\label{DetEqLatticeSquareBarrier}
w\in\mathcal I_\alpha^B(T)
\quad\Longrightarrow\quad
\pi_1(t)\leqslant\pi_2(w)\text{ for some }t\in T.
\end{equation}
Indeed, write $w=q(t_1,\ldots,t_k)$ with
$q(\alpha,\ldots,\alpha)=t\in T$. Compare the first coordinate of
this anchor with the second coordinate of $w$. All variable inputs
increase from $d$ to $\pi_2(t_i)$, and the first coordinate of each
polynomial parameter is at most its second. Monotonicity proves
\eqref{DetEqLatticeSquareBarrier}.

If $L$ contains $M_3$, retain the notation $o,t,a,b,c$ from the
preceding proof and take
\[
\alpha=(a,t),\qquad
T=\{(a,a),(a,t),(b,t)\}.
\]
The first induction contains $(o,t)$, by meeting $(a,t)$ and $(b,t)$.
The polynomial
\[
p(z)=(z\wedge(o,b))\vee(z\wedge(o,c))
\]
has $p(\alpha)=(o,t)$ and $p(a,a)=(o,o)$, so the second induction
contains $(o,o)$. The first does not: neither $a$ nor $b$ is at most
$o$, contradicting \eqref{DetEqLatticeSquareBarrier}.

If $L$ contains $N_5$, with $o<a<b<t$ and remaining element $c$, take
\[
\alpha=(a,t),\qquad
T=\{(b,b),(b,t),(c,t)\}.
\]
The polynomial $q(z)=(z\vee(c,t))\wedge(b,t)$ has
$q(\alpha)=(b,t)\in T$ and $q(c,t)=(o,t)$, so $(o,t)$ belongs to
the first induction. Now
\[
p(z)=(z\wedge(o,a))\vee(z\wedge(o,c))
\]
has $p(\alpha)=(o,t)$ and $p(b,b)=(o,a)$. Thus $(o,a)$ belongs to
the second induction, but not to the first, since neither $b$ nor
$c$ is at most $a$. Again \eqref{DetEqLatticeSquareBarrier} applies.
The forbidden-sublattice criterion therefore proves that
$r_{\mathcal I}^{\mathrm{univ}}(V)\geqslant2$ whenever $V$ is not
distributive. The bounds $o,t$ here are local to the sublattice;
the argument takes place in the full algebra $B$, so no assumption
about named constants in these sublattices is needed.

Finally, in a non-trivial lattice choose $\ell<h$. In its square,
induction at $(\ell,\ell)$ from
$\{(\ell,\ell),(h,h)\}$ adds $(\ell,h)$ by meeting with this
parameter. Hence the universal inductive rank is not zero.
Deduction from a singleton different from the basepoint adds that
basepoint, so the universal deductive rank is not zero either.
\end{proof}

\end{document}
